\begin{helper}
Fix a length-\(k\) star-product tuple \(v\) and the concrete marking rule
of \(\noderef{StarProductConcreteMarked}\).  Let \(U_\ell(i)\) be the
prefix rank-at-most sizes from \(\noderef{StarProductPathRankAtMostSize}\).
Suppose the unmarked steps of \(v\) are assigned labels
\(\psi(i)\in\{0,\ldots,t\}\), and suppose that the sequences
\(U_\ell(i)\) satisfy the hypotheses of \(\noderef{StarProductPathUnmarkedBound}\):
initially \(U_\ell(0)\le N\); at every unmarked step \(i+1\) the labelled
sequence shrinks by a factor \(\rho\); all other labelled sequences are
nonincreasing; the predecessor set for every unmarked step is nonempty; and
the collapse inequality
\[
  N\rho^c<1
\]
holds for every \(c\ge w/(t+1)-1\) in the truncated natural-number sense.
Then the marked-tuple signature of \(v\) has weight at most \(w\).
\end{helper}
\begin{proof}
Let \(\operatorname{sig}\) be the marked-tuple signature of \(v\) from
\(\noderef{StarProductMarkedTupleSignature}\).  By its definition, a
one-bit of \(\operatorname{sig}\) is exactly an unmarked extension step.
Index the step corresponding to \(j\in\{0,\ldots,k-1\}\) by the natural
number \(j+1\).  This index shift is a bijection from the one-bits of
\(\operatorname{sig}\) to the unmarked indices in \(\{1,\ldots,k\}\).
Therefore the binary weight
\[
  |\operatorname{sig}|
\]
is exactly the cardinality of the unmarked-index set counted by
\(\noderef{StarProductPathUnmarkedBound}\).

The hypotheses of the present helper are precisely the initial bound,
shrink bound, monotonicity bound, nonemptiness bound, and collapse bound
required by \(\noderef{StarProductPathUnmarkedBound}\), with
\[
  U_\ell(i)=|U_\ell^{v_{\le i}}|
\]
as defined in \(\noderef{StarProductPathRankAtMostSize}\).  Applying
\(\noderef{StarProductPathUnmarkedBound}\) gives that the shifted unmarked
index set has size at most \(w\).  By the bijective index shift above, this
is exactly
\[
  |\operatorname{sig}|\le w.
\]
\end{proof}
